\documentclass[aspectratio=169]{beamer}
\usetheme{Madrid}
\usecolortheme{seahorse}
\usepackage{amsmath, amssymb}
\usepackage{listings}
\usepackage{xcolor}
\usepackage{tikz}
\usetikzlibrary{arrows.meta, positioning, shapes.geometric, fit, backgrounds}

\lstset{
  basicstyle=\ttfamily\small,
  commentstyle=\color{gray!70}\itshape,
  numbers=left,
  numberstyle=\tiny\color{gray},
  breaklines=true,
  frame=single,
  backgroundcolor=\color{gray!8},
  keepspaces=true,
  showstringspaces=false,
}

% Include a benchmark chart if it has been generated, else show a placeholder so
% the deck still compiles before the scripts have been run.
\newcommand{\benchchart}[1]{%
  \IfFileExists{#1}{%
    \includegraphics[width=\textwidth, height=0.80\textheight, keepaspectratio]{#1}%
  }{%
    \begin{center}\fbox{\parbox{0.8\textwidth}{\centering\vspace{1em}
    \emph{Chart placeholder}\\[0.4em]
    {\footnotesize\texttt{\detokenize{#1}}}\\[0.6em]
    {\footnotesize Run \texttt{bench\_normalization\_charts.py} and
    \texttt{plot\_normalization\_bar\_charts.py} to generate it.}
    \vspace{1em}}}\end{center}%
  }%
}

\title{Normalization Kernels in Rust}
\subtitle{Row-Parallel L1/L2/Max Normalization with In-Place, Zero-Allocation Kernels}
\author{Armin Ebrahimi Saba}
\date{\today}

% ─────────────────────────────────────────────────────────────────────────────

\begin{document}

\begin{frame}
  \titlepage
\end{frame}

\begin{frame}{Outline}
  \tableofcontents
\end{frame}

% ─────────────────────────────────────────────────────────────────────────────
\section{Objective and Kernels}

\begin{frame}{What These Kernels Do}
  \begin{block}{Goal}
    Drop-in Rust replacements for scikit-learn's row-wise \texttt{normalize},
    operating on dense \texttt{f32} matrices, parallelised with Rayon.
  \end{block}
  \vspace{0.4em}
  \textbf{Row-wise vector norms} — each row is scaled independently to unit norm:
  \begin{itemize}
    \item \textbf{L2}: divide row by its Euclidean norm $\|x\|_2$
    \item \textbf{L1}: divide row by its absolute-value sum $\|x\|_1$
    \item \textbf{Max}: divide row by its maximum absolute value $\max_j|x_j|$
  \end{itemize}
  \vspace{0.5em}
  Every kernel comes in a \textbf{copy} form (returns a new array) and an
  \textbf{in-place} form (overwrites the caller's buffer). A zero row is passed
  through unchanged (divide-by-zero guard).
\end{frame}

\begin{frame}{Kernels and Their sklearn Equivalents}
  \begin{center}\small
  \begin{tabular}{lll}
    \hline
    \textbf{Family} & \textbf{Rust kernel(s)} & \textbf{sklearn equivalent} \\
    \hline
    L1  & \texttt{normalize\_l1(\_inplace)}  & \texttt{normalize(X, norm='l1')} \\
    L2  & \texttt{normalize\_l2(\_inplace)}  & \texttt{normalize(X, norm='l2')} \\
    Max & \texttt{normalize\_max(\_inplace)} & \texttt{normalize(X, norm='max')} \\
    \hline
  \end{tabular}
  \end{center}
  \vspace{0.4em}
  \begin{itemize}
    \item All inputs are \texttt{float32}, C-contiguous (row-major)
    \item sklearn keeps \texttt{float32} for these routines (no forced \texttt{f64} upcast)
    \item Data includes $\sim$1\% all-zero rows to exercise the divide-by-zero guard
  \end{itemize}
\end{frame}

\begin{frame}{The Math}
  \textbf{Row norms} (row $x$, output $\hat{x}$):
  \[
    \text{L2:}\ \ \hat{x} = \frac{x}{\|x\|_2}
    \qquad
    \text{L1:}\ \ \hat{x} = \frac{x}{\|x\|_1}
    \qquad
    \text{Max:}\ \ \hat{x} = \frac{x}{\max_j |x_j|}
  \]
  with a guard: a zero row (norm $=0$) is passed through unchanged.
  \vspace{0.8em}
  \begin{alertblock}{Key property — almost no arithmetic per element}
    Computing the norm is a single reduction over the row; applying it is one
    multiply per element. These kernels are \textbf{memory-bandwidth bound}, not
    compute bound --- which shapes every optimisation and every benchmark trend
    that follows.
  \end{alertblock}
\end{frame}

% ─────────────────────────────────────────────────────────────────────────────
\section{Optimisations}

\begin{frame}[fragile]{\#1 — Row-Parallel Rayon}
  Row norms are \textbf{embarrassingly parallel}: each output row depends only on
  the matching input row. We iterate rows and let Rayon work-steal across cores.
\begin{lstlisting}
data.axis_iter_mut(Axis(0))          // one item per row
    .into_par_iter()
    .for_each(|mut row| {
        let norm_sq: f32 = row.iter().map(|x| x * x).sum();
        if norm_sq > 0.0 {
            let inv = 1.0 / norm_sq.sqrt();
            row.mapv_inplace(|x| x * inv);
        }
    });
\end{lstlisting}
  \begin{itemize}
    \item Each row is contiguous (C-order) $\Rightarrow$ unit-stride, full 128 B
          cache lines, prefetcher-friendly
    \item No shared mutable state within a pass $\Rightarrow$ zero synchronisation
    \item All work routed through the \texttt{stratum} global pool
          (\texttt{SKRUB\_RUST\_THREADS}), so it never fights NumPy's BLAS threads
  \end{itemize}
\end{frame}

\begin{frame}[fragile]{\#2 — Single-Pass Copy: Write the Output Exactly Once}
  The copy-returning variant must allocate an output. The naïve pattern
  (\texttt{to\_owned()} then normalise) reads the input, writes a duplicate, then
  overwrites every element again --- \textbf{4N} bytes of RAM traffic.
\begin{lstlisting}
let mut out = Array2::<f32>::uninit((n_rows, n_cols)); // no zero-fill
out.axis_iter_mut(Axis(0)).into_par_iter()
   .zip(src.axis_iter(Axis(0)))
   .for_each(|(mut out_row, in_row)| {
       let inv = 1.0 / l2_norm(in_row);
       for (o, &x) in out_row.iter_mut().zip(in_row.iter()) {
           o.write(x * inv);           // each element written once
       }
   });
\end{lstlisting}
  \begin{itemize}
    \item Uninitialised output, written exactly once $\Rightarrow$ \textbf{2N} traffic
          (1 read + 1 write), no zero-fill pass
    \item Still bound by allocating + first-touching the new buffer (page faults)
  \end{itemize}
\end{frame}

\begin{frame}[fragile]{\#3 — In-Place Kernels: Zero Allocation}
  When the caller does not need the original data, we mutate their NumPy buffer
  directly through a read--write borrow --- \textbf{no allocation, no page faults}.
\begin{lstlisting}
// PyO3 binding takes a read-write view; no copy in, no copy out.
fn normalize_l2_inplace(mut data: PyReadwriteArray2<f32>) {
    let mut view = data.as_array_mut();
    py.allow_threads(|| normalize::normalize_l2_inplace(&mut view));
}
\end{lstlisting}
  \begin{itemize}
    \item The copy variant's cost is dominated by allocating and first-touching a
          fresh output; the in-place form skips that entirely
    \item Leaves only the (parallel, bandwidth-bound) read + write of the data
    \item This is where the \textbf{largest speedups} over sklearn appear ---
          sklearn's \texttt{copy=False} path still does more bookkeeping
    \item \texttt{allow\_threads} releases the GIL so Rayon threads run freely
  \end{itemize}
\end{frame}

\begin{frame}{Memory Layout and Iteration Order}
  \begin{columns}[T]
    \column{0.52\textwidth}
    \textbf{Row-major (C-contiguous) throughout}
    \begin{itemize}
      \item $X$ is used exactly as NumPy stores it --- row-major \texttt{f32},
            \textbf{no transpose, no copy}
      \item Kernels iterate over \texttt{Axis(0)} (rows); each row is a
            contiguous, unit-stride span
      \item One Rayon task per row block --- disjoint slices, no synchronisation
    \end{itemize}
    \textbf{Why row-major is ideal here}
    \begin{itemize}
      \item Row norms are naturally per-row, so row-major storage matches the
            access pattern exactly --- sequential, no strided traversal
    \end{itemize}
    \column{0.48\textwidth}
    \textbf{Consequences}
    \begin{itemize}
      \item 128 B cache line $=$ 32 \texttt{f32}; unit stride uses every byte
      \item Hardware prefetcher stays locked onto the linear scan
      \item Parallelism is over rows, so load balances naturally across cores
    \end{itemize}
    \vspace{0.4em}
    \begin{block}{One-line summary}
      Every kernel is a single (or double) linear sweep of $X$ at unit stride,
      split across cores by row.
    \end{block}
  \end{columns}
\end{frame}

% ─────────────────────────────────────────────────────────────────────────────
\section{Optimisation History}

\begin{frame}{Commit-by-Commit Timeline}
  The commits that shaped the L1/L2/max kernels, oldest to newest:
  \vspace{0.4em}
  \begin{center}\small
  \begin{tabular}{lp{5.4cm}p{4.6cm}}
    \hline
    \textbf{Commit} & \textbf{Change} & \textbf{Bottleneck solved} \\
    \hline
    \texttt{3a5a5cb} & Initial row-wise L1/L2/max normalizers & --- (baseline) \\[3pt]
    \texttt{5985dc4} & Route all Rayon work through the \textbf{global thread pool}
      & Uncontrolled parallelism; oversubscription vs NumPy BLAS \\[3pt]
    \texttt{1c8a8f5} & \textbf{In-place} L1/L2/max via \texttt{PyReadwriteArray2}
      & Copy in / copy out across the FFI boundary \\[3pt]
    \texttt{6220f51} & Copy kernels write \textbf{uninitialised} output once;
      release profile (LTO) & \texttt{to\_owned} then overwrite (4N);
      \texttt{zeros} zero-fill; allocation page-faults \\
    \hline
  \end{tabular}
  \end{center}
\end{frame}

\begin{frame}{The Two Highest-Impact Fixes}
  \begin{columns}[T]
    \column{0.5\textwidth}
    \textbf{Single-pass copy: 4N $\rightarrow$ 2N (\texttt{6220f51})}
    \begin{itemize}
      \item Before: \texttt{to\_owned()} then normalise re-read \& rewrote every
            element --- \textbf{4N} bytes of RAM traffic
      \item After: read input, write the normalised result into an
            \textbf{uninitialised} output exactly once --- \textbf{2N}, no
            zero-fill pass
      \item Same commit: release profile
            (\texttt{opt-level=3}, thin LTO, \texttt{codegen-units=1})
    \end{itemize}
    \column{0.5\textwidth}
    \textbf{In-place kernels: zero allocation (\texttt{1c8a8f5})}
    \begin{itemize}
      \item Copy kernels are dominated by allocating and first-touching
            (page-faulting) a fresh output buffer
      \item In-place variants borrow the NumPy buffer
            (\texttt{PyReadwriteArray2}) and mutate it: no allocation, no page
            faults --- only the parallel read + write remain
      \item This is where the \textbf{largest speedups} over sklearn appear
    \end{itemize}
    \vspace{0.3em}
    \textbf{Thread pool (\texttt{5985dc4})}: all Rayon work runs in the
    \texttt{stratum} global pool, so thread count is controllable and never
    oversubscribes against NumPy's BLAS threads.
  \end{columns}
\end{frame}

% ─────────────────────────────────────────────────────────────────────────────
\section{Experimental Setup}

\begin{frame}{System Specifications}
  \begin{columns}[T]
    \column{0.52\textwidth}
    \textbf{Hardware — Apple MacBook Pro (M1, 2020)}
    \begin{itemize}
      \item SoC: \textbf{Apple M1} (5 nm), 8-core CPU
      \item \textbf{4 Firestorm P-cores} ($\leq$3.2 GHz) +
            \textbf{4 Icestorm E-cores} ($\leq$2.06 GHz)
      \item Cache line: \textbf{128 B} (32 \texttt{f32})
      \item P-core: 128 KB L1d, 192 KB L1i, \textbf{12 MB shared L2}
      \item E-core: 64 KB L1d, \textbf{4 MB shared L2}; $\sim$8 MB System-Level
            Cache (no monolithic L3)
      \item \textbf{8 GB} unified LPDDR4X-4266 ($\sim$68 GB/s)
    \end{itemize}
    \column{0.48\textwidth}
    \textbf{Software}
    \begin{itemize}
      \item macOS 26.5.1 (build 25F80)
      \item Python 3.13 · NumPy 2.3.4 · scikit-learn 1.8.0
      \item Rust extension via PyO3 / maturin, \texttt{--release}
            (\texttt{opt-level=3}), Rayon thread pool
    \end{itemize}
    \vspace{0.3em}
    \begin{alertblock}{Bandwidth is the ceiling}
      At $\sim$68 GB/s, a 2 GB matrix takes $\gtrsim$30 ms just to read once ---
      these bandwidth-bound kernels live right at that limit.
    \end{alertblock}
  \end{columns}
\end{frame}

\begin{frame}{Run Configuration}
  \begin{columns}[T]
    \column{0.5\textwidth}
    \textbf{What is measured}
    \begin{itemize}
      \item Each Rust kernel vs.\ its own sklearn equivalent
      \item Thread counts 1, 2, 4, 8 (machine has 8 cores)
      \item Median of 5 reps (3 for the largest); warm-ups discarded;
            \texttt{gc.collect()} before each rep
      \item Each (kernel, size, threads) cell runs in an isolated subprocess so
            the thread pool initialises with the right count
    \end{itemize}
    \textbf{In-place timing}
    \begin{itemize}
      \item Buffer rebuilt \emph{before} each rep, so the copy cost stays
            \emph{outside} the timed region --- both Rust and sklearn measured the
            same way
    \end{itemize}
    \textbf{Correctness (default on, \texttt{--no-verify} to skip)}
    \begin{itemize}
      \item Before timing, each kernel's Rust output is compared element-wise to
            sklearn's (\texttt{rtol=1e-3}, \texttt{atol=1e-4})
      \item All kernels pass; max abs error $\sim$10$^{-7}$ --- the \texttt{f32}
            row norms agree with sklearn to near machine epsilon
    \end{itemize}
    \column{0.5\textwidth}
    \textbf{Data (synthetic, seed $=$ 42)}
    \begin{itemize}
      \item $X \sim \mathcal{N}(0,1)^{n\times p}$, \texttt{float32}
      \item $\sim$1\% of rows set to all-zero (edge-case stress)
    \end{itemize}
    \vspace{0.4em}
    \textbf{Size tiers}
    \begin{center}\small
    \begin{tabular}{ll}
      \hline
      Tier & Shapes \\
      \hline
      small   & 10k–1M rows, 50–1000 cols \\
      large   & 1M rows, 100 / 500 cols \\
      xlarge  & 2M / 5M rows \\
      xxlarge & 10M$\times$50 ($\sim$2 GB) \\
      \hline
    \end{tabular}
    \end{center}
    \vspace{0.3em}
    \textbf{Reproducible} via \texttt{bench\_normalization\_charts.py}; all
    timings cached in \texttt{normalization\_full\_cache.json}.
  \end{columns}
\end{frame}

% ─────────────────────────────────────────────────────────────────────────────
\section{Benchmark Results — Bar Charts (per size)}

\begin{frame}{sklearn vs Rust — 10k$\times$100 (4 MB)}
  \benchchart{benchmarks/results/normalization_bars_10000x100.png}
  \vspace{-0.5em}
  {\footnotesize Red $=$ sklearn, blue $=$ Rust · log $y$-axis ·
   label above each Rust bar $=$ speedup over sklearn · [small tier]}
\end{frame}

\begin{frame}{sklearn vs Rust — 100k$\times$100 (40 MB)}
  \benchchart{benchmarks/results/normalization_bars_100000x100.png}
  \vspace{-0.5em}
  {\footnotesize In-place kernels lead; copy kernels sit near sklearn · [small tier]}
\end{frame}

\begin{frame}{sklearn vs Rust — 100k$\times$1000 (400 MB)}
  \benchchart{benchmarks/results/normalization_bars_100000x1000.png}
  \vspace{-0.5em}
  {\footnotesize Wide matrix (1000 cols) · [small tier]}
\end{frame}

\begin{frame}{sklearn vs Rust — 1M$\times$50 (200 MB)}
  \benchchart{benchmarks/results/normalization_bars_1000000x50.png}
  \vspace{-0.5em}
  {\footnotesize Tall, narrow matrix · [small tier]}
\end{frame}

\begin{frame}{sklearn vs Rust — 1M$\times$100 (400 MB)}
  \benchchart{benchmarks/results/normalization_bars_1000000x100.png}
  \vspace{-0.5em}
  {\footnotesize [large tier]}
\end{frame}

\begin{frame}{sklearn vs Rust — 1M$\times$500 (2 GB)}
  \benchchart{benchmarks/results/normalization_bars_1000000x500.png}
  \vspace{-0.5em}
  {\footnotesize [large tier]}
\end{frame}

\begin{frame}{sklearn vs Rust — 2M$\times$100 (800 MB)}
  \benchchart{benchmarks/results/normalization_bars_2000000x100.png}
  \vspace{-0.5em}
  {\footnotesize Streaming straight from DRAM at the bandwidth limit · [xlarge tier]}
\end{frame}

\begin{frame}{sklearn vs Rust — 5M$\times$50 (1 GB)}
  \benchchart{benchmarks/results/normalization_bars_5000000x50.png}
  \vspace{-0.5em}
  {\footnotesize [xlarge tier]}
\end{frame}

\begin{frame}{sklearn vs Rust — 10M$\times$50 (2 GB)}
  \benchchart{benchmarks/results/normalization_bars_10000000x50.png}
  \vspace{-0.5em}
  {\footnotesize A single pass is $\gtrsim$30 ms of pure memory traffic · [xxlarge tier]}
\end{frame}

% ─────────────────────────────────────────────────────────────────────────────
\section{Benchmark Results — Speedup vs Size (per thread count)}

\begin{frame}{Speedup over sklearn — 1 thread}
  \benchchart{benchmarks/results/normalization_groupA_speedup_t1.pdf}
  \vspace{-0.5em}
  {\footnotesize x $=$ input size · one line per kernel · in-place kernels dashed}
\end{frame}

\begin{frame}{Speedup over sklearn — 2 threads}
  \benchchart{benchmarks/results/normalization_groupA_speedup_t2.pdf}
\end{frame}

\begin{frame}{Speedup over sklearn — 4 threads}
  \benchchart{benchmarks/results/normalization_groupA_speedup_t4.pdf}
\end{frame}

\begin{frame}{Speedup over sklearn — 8 threads}
  \benchchart{benchmarks/results/normalization_groupA_speedup_t8.pdf}
\end{frame}

% ─────────────────────────────────────────────────────────────────────────────
\section{Benchmark Results — Thread Scaling (per size)}

\begin{frame}{Thread Scaling — 10k$\times$100 (small)}
  \benchchart{benchmarks/results/normalization_groupB_speedup_10000x100.pdf}
  \vspace{-0.5em}
  {\footnotesize x $=$ thread count · one line per kernel · speedup over sklearn}
\end{frame}

\begin{frame}{Thread Scaling — 1M$\times$500 (large)}
  \benchchart{benchmarks/results/normalization_groupB_speedup_1000000x500.pdf}
\end{frame}

\begin{frame}{Thread Scaling — 2M$\times$100 (xlarge)}
  \benchchart{benchmarks/results/normalization_groupB_speedup_2000000x100.pdf}
\end{frame}

\begin{frame}{Thread Scaling — 10M$\times$50 (xxlarge)}
  \benchchart{benchmarks/results/normalization_groupB_speedup_10000000x50.pdf}
  \vspace{-0.5em}
  {\footnotesize The remaining sizes and all runtime-metric charts exist as
   individual files under \texttt{benchmarks/results/}}
\end{frame}

% ─────────────────────────────────────────────────────────────────────────────
\section{Why These Trends}

\begin{frame}{Why Copy Kernels Only Tie sklearn, but In-Place Wins}
  \begin{columns}[T]
    \column{0.5\textwidth}
    \textbf{Copy variants: near parity at low threads}
    \begin{itemize}
      \item sklearn's \texttt{normalize} is already a tight, vectorised C loop ---
            there is almost no arithmetic to out-optimise
      \item Our copy kernel must \textbf{allocate} the output and pay
            \textbf{first-touch page faults} on every fresh page
      \item At 1 thread the allocation + page-fault cost can leave Rust
            \emph{slightly behind} sklearn ($\sim$0.3--0.4$\times$)
    \end{itemize}
    \column{0.5\textwidth}
    \textbf{In-place variants: the real win}
    \begin{itemize}
      \item No allocation, no page faults --- only the parallel read + write
      \item Reaches $\sim$3--4$\times$ at 4 threads on large inputs
      \item This is the honest source of the advantage: \textbf{memory-traffic and
            allocation savings}, amplified by threading --- \emph{not} raw compute
    \end{itemize}
    \vspace{0.3em}
    \textbf{Takeaway}: for bandwidth-bound ops, \emph{avoiding memory work} beats
    \emph{doing arithmetic faster}.
  \end{columns}
\end{frame}

\begin{frame}{Why Speedup Grows With Data Size, Then Saturates}
  \begin{columns}[T]
    \column{0.5\textwidth}
    \textbf{Small inputs: fixed costs dominate}
    \begin{itemize}
      \item Python$\leftrightarrow$Rust FFI, thread-pool wake-up and Rayon
            scheduling are constant per call
      \item On a 4 MB matrix they are a large fraction of the total, so speedups
            are muted
    \end{itemize}
    \textbf{Growing size: overheads amortise}
    \begin{itemize}
      \item More elements per call $\Rightarrow$ the fixed costs shrink relative to
            useful work, and parallel efficiency rises
    \end{itemize}
    \column{0.5\textwidth}
    \textbf{Large inputs: the bandwidth wall}
    \begin{itemize}
      \item Once $X$ exceeds cache, the kernel streams from DRAM at
            $\sim$68 GB/s --- a hard ceiling
      \item Extra threads beyond the memory channels ($\sim$4 P-cores) add little:
            the bus, not the ALUs, is the bottleneck
      \item So thread scaling \textbf{plateaus at 4T}, and 8T (E-cores) rarely
            helps for these kernels
    \end{itemize}
    \vspace{0.3em}
    \textbf{Wide vs tall} matters little: the work is proportional to the total
    element count, and every row block is an independent, unit-stride sweep.
  \end{columns}
\end{frame}

% ─────────────────────────────────────────────────────────────────────────────
\section{Summary}

\begin{frame}{Summary}
  \begin{columns}[T]
    \column{0.5\textwidth}
    \textbf{What makes the kernels fast}
    \begin{itemize}
      \item \textbf{Row-parallel Rayon} over the global thread pool
      \item \textbf{Single-pass copy}: uninitialised output written once (2N vs 4N)
      \item \textbf{In-place kernels}: zero allocation, no page faults
      \item \textbf{Row-major, unit-stride} access throughout; zero-copy FFI
    \end{itemize}
    \column{0.5\textwidth}
    \textbf{What the benchmarks show}
    \begin{itemize}
      \item These ops are \textbf{memory-bandwidth bound} --- sklearn is already
            well vectorised
      \item Copy kernels $\approx$ sklearn; \textbf{in-place kernels win big}
            ($\sim$3--4$\times$ at 4T)
      \item Speedup grows with size until the DRAM bandwidth wall
      \item Scaling plateaus at $\sim$4 threads (memory channels, not cores)
    \end{itemize}
    \vspace{0.4em}
    \textbf{Guidance}: prefer in-place kernels and 4 threads; extra cores buy
    little once bandwidth-bound.
  \end{columns}
\end{frame}

\end{document}
